\begin{tabular}{@{}>{\raggedright\arraybackslash}p{0.22\linewidth}>{\raggedright\arraybackslash}p{0.28\linewidth}>{\raggedright\arraybackslash}p{0.43\linewidth}@{}}
\toprule
Phenomenon (A) & Measure (B) & Interpretation in tennis (C) \\
\midrule
Who wins, and who holds prestige & In- and out-degree, PageRank, HITS authority and hub scores & In-degree counts distinct opponents defeated; out-degree counts distinct opponents lost to. PageRank weights wins by opponent prestige; an authority defeats strong contenders, while a hub loses frequently to top authorities by advancing deep into draws. \\
\addlinespace[3pt]
How concentrated prestige and wins are & Share of PageRank held by the top 10, Gini coefficient of in-strength & A high top-10 share concentrates prestige. A high Gini means wins are unevenly distributed. \\
\addlinespace[3pt]
Core players and the wider field & $k$-core decomposition, Borgatti-Everett core-periphery model & A $k$-core keeps players with at least $k$ ties inside the set. The fitted model compares within-core and between-group densities; neither requires every core pair to have played. \\
\addlinespace[3pt]
How connected the tour is & Components, average distance, diameter, global efficiency, ratio to a random reference & Components count separate groups; distances count steps between players. Short paths and high efficiency indicate a more connected field. \\
\addlinespace[3pt]
Dependence on top players & Largest component after removing players at random and in decreasing PageRank order & A faster loss of connectivity after removing high-PageRank players shows greater reliance on leading players. \\
\addlinespace[3pt]
Closed rivalries and win hierarchy & Clustering, transitivity, triad census & Clustering measures closed opponent ties locally; transitivity measures closed triples overall. Triads distinguish transitive results from cycles of wins. \\
\addlinespace[3pt]
Whether opponent counts have a heavy tail & Power-law fit of the degree distribution and likelihood-ratio tests against lognormal, exponential and truncated power law & A power-law tail would mean a few players face many more opponents. Tournament draws and season length set a practical upper bound, so alternatives are compared as well. \\
\addlinespace[3pt]
Do highly ranked or active players meet each other & Degree and ranking-band assortativity, band mixing against permutation nulls & Positive assortativity means similar-degree players meet. Band-mixing ratios compare observed matches with shuffled ranks, including a shuffle that controls opponent counts. \\
\addlinespace[3pt]
Whether players occupy similar positions & Cosine similarity of opponent vectors, QAP correlation with the ranking gap & Cosine similarity compares opponent profiles. QAP tests whether profile similarity tracks the gap in official ranking. \\
\addlinespace[3pt]
Structure beyond opponent counts & Configuration model, $z$-scores and permutation $p$-values, small-world index $\sigma$ & Degree-preserving nulls keep each player's opponent count. Differences from them show patterns not explained by those counts alone. \\
\addlinespace[3pt]
Groups and career cohorts & Louvain communities and modularity & Modularity measures how strongly matches fall within detected groups. Here, groups can reflect players whose careers overlapped in time. \\
\bottomrule
\end{tabular}
